% =============================================================================
% test-abnt-refs-citacoes.tex
% Documento de teste para abnt-refs.sty e abnt-citacoes.sty
%
% Demonstra referências (NBR 6023:2025) e citações (NBR 10520:2023).
%
% Compilação (requer Biber instalado):
%   TEXINPUTS=./src: latexmk -pdf -output-directory=tests tests/test-abnt-refs-citacoes.tex
%
% Para sistema numérico, alterar: \usepackage[sistema=numerico]{abnt-citacoes}
%
% Checklist:
%   1  — Compilação sem erros (XeLaTeX + Biber)
%   2  — Lista de referências formatada (espaçamento, alinhamento)
%   3  — Título REFERÊNCIAS centralizado, maiúsculas, negrito
%   4  — Backref (páginas de citação após cada referência)
%   5  — Citação autor-data inline (\textcite)
%   6  — Citação autor-data parentética (\parencite)
%   7  — Citação com et al. (4+ autores)
%   8  — Citação direta curta (até 3 linhas, entre aspas)
%   9  — Citação direta longa (ambiente citacaodireta)
%   10 — Citação de citação (apud)
%   11 — Sistema numérico (trocar opção para testar)
%   12 — Tipos de referência variados (livro, artigo, tese, legislação, site)
%   13 — Compatibilidade hyperref
%   14 — Compatibilidade abnt.cls
% =============================================================================

\documentclass[a4paper, 12pt]{report}
\usepackage[T1]{fontenc}
\usepackage[brazil]{babel}
\usepackage{abnt-numeracao}
\usepackage{abnt-sumario}
\usepackage{abnt-citacoes}            % carrega abnt-refs automaticamente
\usepackage{hyperref}                  % checklist 13
\usepackage{fancyhdr}
\usepackage{setspace}
\onehalfspacing
\pagestyle{fancy}
\fancyhf{}
\fancyhead[R]{\thepage}
\renewcommand{\headrulewidth}{0pt}

\addbibresource{refs.bib}

\begin{document}

% =====================================================================
% PRÉ-TEXTUAIS
% =====================================================================
\pagestyle{empty}

\begin{titlepage}
  \centering\vfill
  {\Large\bfseries TESTE DE REFERÊNCIAS E CITAÇÕES}\par
  \vfill
  Cidade\par 2025
\end{titlepage}

% =====================================================================
% SUMÁRIO
% =====================================================================
\imprimirsumario

% =====================================================================
% ELEMENTOS TEXTUAIS
% =====================================================================

% =====================================================================
% CAPÍTULO 1 — Citações autor-data (checklist 5, 6)
% =====================================================================
\chapter{Citações no sistema autor-data}

% Checklist 5: citação inline (\textcite)
Segundo \textcite{severino2007}, a metodologia do trabalho científico é
fundamental para a formação acadêmica.

% Checklist 6: citação parentética (\parencite)
A metodologia científica é considerada essencial na formação universitária
\parencite{severino2007}.

% Checklist 7: citação com et al. (4+ autores)
As reformas econômicas no Brasil foram amplamente discutidas na literatura
\parencite{urani1994}.

% Artigo de periódico
A análise de políticas públicas educacionais revela desafios estruturais no
sistema brasileiro \parencite{souza2019}.

% Tese
Estudos sobre impactos ambientais em áreas urbanas demonstram a necessidade de
planejamento territorial \parencite{silva2015}.

% =====================================================================
% CAPÍTULO 2 — Citações diretas (checklist 8, 9)
% =====================================================================
\chapter{Citações diretas}

% Checklist 8: citação direta curta (até 3 linhas, entre aspas)
\section{Citação direta curta}

Conforme \textcite[p.~23]{severino2007}, ``a metodologia é o instrumento que
permite ao pesquisador organizar sua investigação de forma sistemática e
rigorosa.''

% Checklist 9: citação direta longa (ambiente citacaodireta)
\section{Citação direta longa}

A importância da metodologia na pesquisa científica é assim descrita:

\begin{citacaodireta}
A pesquisa científica exige do investigador uma postura metodológica que vai
além da simples aplicação de técnicas. Trata-se de uma atitude intelectual
que envolve questionamento constante, rigor na coleta e análise dos dados,
e capacidade de articular teoria e prática de forma coerente e fundamentada.
O pesquisador deve estar preparado para revisar suas hipóteses e adaptar
seus métodos conforme os resultados parciais o indiquem
\parencite[p.~45--46]{severino2007}.
\end{citacaodireta}

O trecho acima demonstra o uso do ambiente \texttt{citacaodireta}, com recuo
de 4~cm, fonte menor e espaçamento simples.

% =====================================================================
% CAPÍTULO 3 — Citação de citação e outros casos (checklist 10)
% =====================================================================
\chapter{Casos especiais}

% Checklist 10: citação de citação (apud)
\section{Citação de citação}

\textcite{piaget1975} apud \textcite{moreira2011} defende que o
desenvolvimento cognitivo ocorre por meio de processos de assimilação e
acomodação.

% Desambiguação por letra (mesmo autor, mesmo ano)
\section{Desambiguação por ano}

A educação inclusiva foi discutida em dois estudos distintos
\parencite{santos2020a,santos2020b}.

% Múltiplas citações
\section{Múltiplas citações}

Diversos autores trataram de metodologia de pesquisa
\parencite{severino2007,demo2003,silva2015}.

% =====================================================================
% CAPÍTULO 4 — Tipos variados de referência (checklist 12)
% =====================================================================
\chapter{Referências variadas}

% Legislação
A legislação educacional brasileira é complexa \parencite{brasil2014}.

% Site / recurso eletrônico
Os dados populacionais mais recentes estão disponíveis no portal do censo
\parencite{ibge2023}.

% Trabalho em evento
Técnicas computacionais têm sido aplicadas na classificação de solos
\parencite{oliveira2018}.

% Capítulo de livro
A pesquisa como princípio educativo é discutida por \textcite{demo2003}.

% =====================================================================
% PÓS-TEXTUAIS — checklist 2, 3, 4
% =====================================================================
\imprimirreferencias

\end{document}
